\chapter{Tương lai}
\markboth{Tương lai}{The Future}

\section{Sau này không ai trả lương cho sự lười biếng của em.}
\begin{truyen}{Sau này không ai trả lương cho sự lười biếng của em.}{Wake up.}
\chuhoa{C}{âu nói khiến cả lớp cười nhẹ rồi im. Đi học còn có người nhắc, có thầy cô kéo; ra đời rồi, sự lười chỉ làm mình bị loại khỏi cơ hội chứ không được thưởng công.}
\textit{(The line made the class laugh lightly and then fall silent. In school people still remind and pull you along; in real life laziness only removes you from opportunity, it is never rewarded.)}
\end{truyen}

\begin{baihoc}
Những gì em dung túng hôm nay có thể trở thành điểm yếu nghề nghiệp của ngày mai.
\textit{(What you tolerate today may become a professional weakness tomorrow.)}
\end{baihoc}

\begin{ghinhoanh}
\textbf{Short English to remember:}

Laziness earns nothing in the real world.

\end{ghinhoanh}

\begin{tuhocnhanh}
1. Thói quen nào của em sẽ gây hại khi đi làm sau này? \textit{Which habit of yours may hurt you later at work?}

2. Em có thể sửa từ bây giờ không? \textit{Can you fix it from now on?}
\end{tuhocnhanh}
\ngancach

\section{Cơ hội không ghé nhà người chỉ biết nói để mai.}
\begin{truyen}{Cơ hội không ghé nhà người chỉ biết nói để mai.}{Wake up.}
\chuhoa{C}{âu này không dài nhưng đụng đúng căn bệnh trì hoãn. Cơ hội thường đi qua người biết chuẩn bị và biết bắt đầu, chứ không đợi người cứ xin thêm một ngày nữa.}
\textit{(The line was short but hit directly at procrastination. Opportunity usually passes by those who prepare and begin, not those who keep asking for one more day.)}
\end{truyen}

\begin{baihoc}
Muốn gặp cơ hội, em phải trở thành người sẵn sàng trước đã.
\textit{(If you want to meet opportunity, you must first become ready for it.)}
\end{baihoc}

\begin{ghinhoanh}
\textbf{Short English to remember:}

Opportunity visits the prepared, not the postponing.

\end{ghinhoanh}

\begin{tuhocnhanh}
1. Em đang hoãn điều gì mà có thể làm lỡ cơ hội? \textit{What are you postponing that may cost you an opportunity?}

2. Em cần chuẩn bị gì ngay từ bây giờ? \textit{What should you prepare right now?}
\end{tuhocnhanh}
\ngancach

\section{Điểm kém chỉ là cảnh báo; hậu quả thật nằm ở năng lực rỗng.}
\begin{truyen}{Điểm kém chỉ là cảnh báo; hậu quả thật nằm ở năng lực rỗng.}{Wake up.}
\chuhoa{T}{hầy muốn học sinh hiểu rằng điểm thấp không phải bi kịch lớn nhất. Điều đáng sợ hơn là qua nhiều năm mà năng lực thật vẫn trống, vì khi ấy em bước ra đời với vẻ ngoài có bằng cấp nhưng bên trong không đủ sức làm việc.}
\textit{(The teacher wanted students to understand that low grades are not the greatest tragedy. What is more frightening is leaving school with empty real ability, because then you enter life with credentials outside but not enough strength inside.)}
\end{truyen}

\begin{baihoc}
Điểm số báo động sớm để em kịp sửa trước khi năng lực thật trở nên quá yếu.
\textit{(Grades are an early warning so you can improve before real ability becomes too weak.)}
\end{baihoc}

\begin{ghinhoanh}
\textbf{Short English to remember:}

Low grades warn you; empty ability harms you.

\end{ghinhoanh}

\begin{tuhocnhanh}
1. Em đang sợ điểm hay đang sợ sự rỗng thật sự? \textit{Are you afraid of grades or of real emptiness?}

2. Em cần lấp chỗ trống nào trong năng lực của mình? \textit{Which gap in your ability do you need to fill?}
\end{tuhocnhanh}
\ngancach

\section{Em đang học cho bài kiểm tra hay học để không bị đời kiểm tra đau hơn?}
\begin{truyen}{Em đang học cho bài kiểm tra hay học để không bị đời kiểm tra đau hơn?}{Wake up.}
\chuhoa{C}{âu hỏi làm nhiều em ngẩn ra. Bài kiểm tra ở trường chỉ là bản nhẹ; còn đời sẽ hỏi bằng công việc, áp lực, lựa chọn và trách nhiệm thật.}
\textit{(The question made many students pause. School tests are only the lighter version; life will ask through work, pressure, choices, and real responsibility.)}
\end{truyen}

\begin{baihoc}
Học cho sâu là cách giảm bớt những cú kiểm tra đau của cuộc đời.
\textit{(Learning deeply is one way to reduce the painful tests of life.)}
\end{baihoc}

\begin{ghinhoanh}
\textbf{Short English to remember:}

School tests are not the hardest tests in life.

\end{ghinhoanh}

\begin{tuhocnhanh}
1. Em đang học đối phó hay học lâu dài? \textit{Are you studying to survive or to live well later?}

2. Nếu đời kiểm tra kỹ năng của em hôm nay, em có vững không? \textit{If life tested your skills today, would you stand firm?}
\end{tuhocnhanh}
\ngancach

\section{Nghèo kiến thức còn nguy hơn nghèo tiền.}
\begin{truyen}{Nghèo kiến thức còn nguy hơn nghèo tiền.}{Wake up.}
\chuhoa{T}{hầy không coi nhẹ sự khó khăn vật chất, nhưng thầy nói thẳng: nghèo tiền còn có thể kiếm dần nếu còn năng lực; nghèo kiến thức và kỹ năng thì rất dễ bị mắc kẹt lâu hơn.}
\textit{(The teacher did not belittle material hardship, but he spoke plainly: lacking money can still be improved if ability remains; lacking knowledge and skills can trap a person much longer.)}
\end{truyen}

\begin{baihoc}
Tri thức là một loại vốn giúp em tự đứng lên bền hơn.
\textit{(Knowledge is a kind of capital that helps you stand up more steadily.)}
\end{baihoc}

\begin{ghinhoanh}
\textbf{Short English to remember:}

Knowledge is wealth that keeps creating value.

\end{ghinhoanh}

\begin{tuhocnhanh}
1. Em đang tích lũy vốn tri thức của mình ra sao? \textit{How are you building your knowledge capital?}

2. Vì sao kiến thức có thể giúp em bền hơn tiền? \textit{Why can knowledge support you longer than money?}
\end{tuhocnhanh}
\ngancach

\section{Người chăm hôm nay đỡ phải van xin ngày mai.}
\begin{truyen}{Người chăm hôm nay đỡ phải van xin ngày mai.}{Wake up.}
\chuhoa{C}{âu nói nghe nặng nhưng rất thật. Khi đã chuẩn bị tốt, em có nhiều quyền chọn hơn; còn khi tay trắng về năng lực, người ta thường phải xin xỏ những điều đáng lẽ mình có thể tự tạo ra.}
\textit{(The line sounded heavy but true. When you prepare well, you have more choices; when your ability is empty, you often have to beg for things you could have built yourself.)}
\end{truyen}

\begin{baihoc}
Sự siêng năng hôm nay là một dạng giữ phẩm giá cho ngày mai.
\textit{(Today's diligence is one way of protecting tomorrow's dignity.)}
\end{baihoc}

\begin{ghinhoanh}
\textbf{Short English to remember:}

Work now so you won't beg later.

\end{ghinhoanh}

\begin{tuhocnhanh}
1. Em muốn ngày mai mình ở thế chủ động hay bị động? \textit{Do you want your future self to be active or helpless?}

2. Hôm nay em phải chăm ở đâu để đạt điều đó? \textit{Where must you work hard today to reach that?}
\end{tuhocnhanh}
\ngancach

\section{Nếu em không tự nâng mình lên, đời sẽ xếp em xuống chỗ thấp hơn.}
\begin{truyen}{Nếu em không tự nâng mình lên, đời sẽ xếp em xuống chỗ thấp hơn.}{Wake up.}
\chuhoa{C}{âu này làm nhiều em giật mình vì nó không màu mè. Cuộc sống không giữ chỗ đẹp cho người ngừng cố gắng; ai đứng yên quá lâu sẽ tự bị tụt lại phía sau.}
\textit{(The sentence startled many students because it was not decorative at all. Life does not reserve a good place for those who stop trying; anyone who stands still too long falls behind.)}
\end{truyen}

\begin{baihoc}
Muốn ở vị trí tốt hơn, em phải tự nâng trình độ và ý thức của mình lên trước.
\textit{(If you want a better place, you must first raise your level and your awareness.)}
\end{baihoc}

\begin{ghinhoanh}
\textbf{Short English to remember:}

Raise yourself, or life will place you lower.

\end{ghinhoanh}

\begin{tuhocnhanh}
1. Em đang đứng yên ở mặt nào? \textit{In which area are you standing still?}

2. Em có thể nâng mình lên bằng việc gì? \textit{How can you raise yourself?}
\end{tuhocnhanh}
\ngancach

\section{Tương lai không thương người biết đúng mà không làm.}
\begin{truyen}{Tương lai không thương người biết đúng mà không làm.}{Wake up.}
\chuhoa{T}{hầy biết nhiều học sinh không thiếu hiểu biết; các em thiếu hành động. Biết điều đúng mà cứ để đó sẽ không cứu được ai, kể cả chính mình.}
\textit{(The teacher knew that many students do not lack understanding; they lack action. Knowing the right thing and leaving it untouched will save no one, not even yourself.)}
\end{truyen}

\begin{baihoc}
Tri thức chỉ có sức mạnh khi nó bước ra thành hành động.
\textit{(Knowledge has power only when it steps out into action.)}
\end{baihoc}

\begin{ghinhoanh}
\textbf{Short English to remember:}

Knowing is not enough if you never act.

\end{ghinhoanh}

\begin{tuhocnhanh}
1. Điều đúng nào em đã biết nhưng chưa làm? \textit{What right thing do you know but have not done?}

2. Điều gì đang cản em hành động? \textit{What is stopping you from acting?}
\end{tuhocnhanh}
\ngancach

\section{Mỗi ngày lười là một viên gạch xây nên hối tiếc.}
\begin{truyen}{Mỗi ngày lười là một viên gạch xây nên hối tiếc.}{Wake up.}
\chuhoa{C}{âu này nghe vừa hình ảnh vừa đáng sợ. Hối tiếc lớn hiếm khi được tạo trong một ngày; nó thường được xây dần từ những ngày mình biết mình nên làm gì nhưng lại bỏ qua.}
\textit{(The line was both vivid and frightening. Great regret is rarely made in a single day; it is usually built slowly from many days when we knew what to do but ignored it.)}
\end{truyen}

\begin{baihoc}
Đừng xây hối tiếc bằng những lựa chọn dễ dãi lặp đi lặp lại.
\textit{(Do not build regret through repeated easy choices.)}
\end{baihoc}

\begin{ghinhoanh}
\textbf{Short English to remember:}

Daily laziness builds future regret.

\end{ghinhoanh}

\begin{tuhocnhanh}
1. Em có đang đặt “một viên gạch hối tiếc” mỗi ngày không? \textit{Are you laying “a brick of regret” each day?}

2. Hôm nay em có thể tháo viên nào ra? \textit{Which one can you remove today?}
\end{tuhocnhanh}
\ngancach

\section{Cái giá của lười học thường đến muộn nên nhiều người chủ quan.}
\begin{truyen}{Cái giá của lười học thường đến muộn nên nhiều người chủ quan.}{Wake up.}
\chuhoa{T}{hầy bảo: ``Chính vì cái giá không đến ngay nên người ta tưởng không sao. Nhưng khi nó đến~--- bằng sự thiếu tự tin, thiếu kỹ năng, thiếu cơ hội~--- thì lại khó sửa hơn rất nhiều.''}
\textit{(The teacher said: ``Because the price does not arrive immediately, people think everything is fine. But when it finally comes—in the form of low confidence, weak skills, and fewer opportunities—it becomes much harder to fix.'')}
\end{truyen}

\begin{baihoc}
Đừng để sự chậm trễ của hậu quả làm em coi nhẹ nguyên nhân.
\textit{(Do not let delayed consequences make you underestimate the cause.)}
\end{baihoc}

\begin{ghinhoanh}
\textbf{Short English to remember:}

Delayed cost is still real cost.

\end{ghinhoanh}

\begin{tuhocnhanh}
1. Em có đang chủ quan vì chưa thấy hậu quả ngay không? \textit{Are you careless because you have not seen the consequence yet?}

2. Nếu nhìn xa hơn, em sẽ đổi điều gì hôm nay? \textit{If you looked further ahead, what would you change today?}
\end{tuhocnhanh}
\ngancach